\section{结论}
\label{sec:conclusion}

RAG 通过外部知识扩展了语言模型的能力边界，但固定的“一次检索—生成”流程同时带来不必要调用、查询失配、证据噪声和复杂问题覆盖不足。自适应 RAG 的演进可以概括为控制点的逐步增加：Adaptive-RAG 在输入阶段选择路径，FLARE 和 DRAGIN 在生成中识别信息需求，HyDE 与 RQ-RAG 改善查询表达，CRAG 对失败检索采取纠正行动，Self-RAG 将检索与反思纳入学习过程，IRCoT 和 Iter-RetGen 则让检索随推理逐步展开。

这些方法共同表明，RAG 不应被理解为给提示词附加文档，而应被理解为在参数知识与外部证据之间进行动态决策。本文进一步认为，其关键科学问题是估计信息缺口及补充证据的价值：问题复杂度不是检索必要性的充分条件，Token 置信度不是事实可靠性的充分条件，多轮反思也不是证据充分性的充分条件。未来研究需要更加重视显式证据状态、可审计停止条件、独立验证及质量—成本联合评价。自适应系统的目标不是执行最多步骤，而是在必要时采取尽可能少且有效的步骤，并在证据不足时知道自己应当停止。
